\section{Kompleksna dinamika}

\subsection{Kompleksne periodične točke}

\noindent Definicijo privlačne in odbojne točke nam v kompleksni ravnini poenostavi naslednji izrek.

\begin{izrek}
    Naj bo \(f\) cela funkcija in \(z^*\) njena fiksna točka. Potem je \(z^*\) privlačna natanko tedaj, ko je \(|f' (z^*)| < 1\). Dodatno je \(z^*\) odbojna natanko tedaj, ko je \(|f' (z^*)| > 1\).
\end{izrek}

\begin{dokaz}
    Naj bo \(z^*\) privlačna. Ker za vsak \(z \in \Delta (z^*, r)\) zaporedje \(f^n (z)\) enakomerno konvergira proti \(z^*\), obstaja \(n_0 \in \NN\), tako da \(|f^n (z) - z^*| < r\) za vsak \(z \in \Delta (z^*, r)\) in \(n \geq n_0\). Opazujemo funkcijo \(\phi (z) = (z - z^*) / r\), ki za vsak \(n \geq n_0\) konjugira \(f^n\) in \(F_n \coloneq \phi \circ f^n \circ \phi^{-1}\). Ko poenostavimo dobimo
    \[F_n (w) = \frac{f^n (r w + z^*) - z^*}{r},\]
    kjer je vsaka \(F_n\) definirana na \(\DD\), ima vrednosti v \(\overline{\DD}\) in zadošča \(F_n (0) = 0\). Schwarzova lema nam pove, da je \(|F_n' (0)| \leq 1\). Recimo, da velja \(|F_n' (0)| = 1\). Spet po Schwarzovi lemi je \(F_n\) rotacija, torej \(F_n (w) = a_n w\) za \(|a_n| = 1\). Tedaj velja
    \[\abs{F_n (w)} = \abs{w} = \frac{\abs{f^n (r w + z^*) - z^*}}{r}\]
    ali ekvivalentno \(\abs{f^n (z) - z^*} = \abs{z - z^*}\) za vsak \(z \in \Delta (z^*, r)\). Ker to pomeni, da \(f^n (z)\) ne konvergira proti \(z^*\), smo prišli do protislovja.

    Obratno, naj obstaja \(k \in \RR\), da vejla \(|f' (z^*)| < k < 1\). Iz definicije odvoda sledi, da obstaja \(r > 0\), da za vsak \(z \in \Delta (z^*, r)\) z \(z \neq z^*\) velja
    \[\abs{\frac{f (z) - f (z^*)}{z - z^*}} < k.\]
    To je ekvivalentno \(|f (z) - z^*| \leq k |z - z^*|\). Induktivno nato pokažemo, da velja \(|f^n (z) - z^*| \leq k^n |z - z^*|\) za \(z \in \Delta (z^*, r)\) in \(z \neq z^*\). Ker je \(\lim_{n \to \infty} k^n = 0\), je \(z^*\) res privlačna.

    Naj bo sedaj \(z^*\) odbojna. Iz zgornjega dokaza vemo, da velja \(|f' (z^*)| \geq 1\). Ker \(|f' (z^*)| \neq 0\), je \(f\) lokalni biholomorfizem na neki okolici točke \(z^*\). Torej lahko izberemo kompaktno okolico \(V\) točke \(z^*\), na kateri je \(f\) biholomorfna, \(z^*\) pa točke v \(V\) odbije, torej njihove orbite divergirajo iz \(V\). Za vsak \(k \in \NN\) nato definiramo
    \[V_k \coloneq V \cap f^{-1} (V) \cap \cdots \cap f^{- k} (V).\]
    Vsaka \(V_k\) je kompaktna okolica \(z^*\), ki vsebuje vse točke \(z\), za katere \(f^j (z) \in V\) za \(1 \leq j \leq k\). Če označimo \(V_0 = V\), je zaporedje \((V_k)_{k \geq 0}\) padajoče zaporedje kompaktnih množic z \(\bigcap_k V_k = \set{z^*}\). Poleg tega iz definicije \(V_k\) sledi zveza \(f (V_k) = V_{k - 1} \cap f (V)\). Ker obstaja \(k_0 \in \NN\), tako da za vsak \(k \geq k_0\) velja \(V_k \subset f (V)\), velja tudi \(f (V_k) = V_{k - 1}\) za \(k \geq k_0\). Izberemo torej \(r > 0\), da \(\Delta (z^*, r) \subset V_{k_0}\) in tak \(h > k_0\), da je \(V_h \subset \Delta (z^*, r / 2)\). Z \(g\) označimo \((h - k_0)\)-to iteracijo \(f^{-1}\), to je \(g = (f^{-1})^{h - k_0}\). Tedaj vidimo, da velja
    \[g \prt{\Delta (z^*, r)} \subset V_h \subset \Delta (z^*, r).\]
    Ker sta \(\Delta (z^*, r)\) in \(g \prt{\Delta (z^*, r)}\) enostavno povezani območji, lahko uporabimo Schwarzovo lemo, ki nam pove, da je \(|g' (z^*)| < 1\). Ker velja \(g' (z^*) = (1 / f' (z^*))^{h - k_0}\), je res \(|f' (z^*)| > 1\).

    Obratno, naj obstaja \(k > 1\), da velja \(|f' (z^*)| > k\). Podobno kot prej obstaja \(r > 0\), da je \(|f (z) - f (z^*)| \geq k |z - z^*|\) za vsak \(z \in \Delta (z^*, r)\) in \(z \neq z^*\). Recimo, da obstaja \(z_0\), da velja \(f^n (z_0) \in \Delta (z^*, r)\) za vsak \(n \in \NN\). Tedaj lahko induktivno pokažemo, da velja \(|f^n (z_0) - z^*| \geq k^n |z_0 - z^*|\). Ker je \(\lim_{n \to \infty} k^n = \infty\), smo prišli do protislovja. Torej je \(z^*\) res odbojna.
\end{dokaz}

\noindent V luči zgornjega izreka lahko definicijo razširimo.

\begin{definicija}
    Naj bo \(U \subseteq \CC\) odprta in \(f \colon U \to \CC\) holomorfna. Naj bo \(w \in U\) fiksna točka preslikave \(f\). Število \(\lambda \coloneq f' (w)\) imenujemo \emph{multiplikator preslikave \(f\) v točki \(w\)}. Fiksna točka \(w\) je
    \begin{itemize}
        \item \emph{superprivlačna}, če je \(\lambda = 0\);
        \item \emph{privlačna}, če je \(|\lambda| < 1\);
        \item \emph{odbojna}, če je \(|\lambda| > 1\);
        \item \emph{racionalno nevtralna}, če je \(|\lambda| = 1\) in \(\lambda^n = 1\) za kak \(n \in \NN\);
        \item \emph{iracionalno nevtralna}, če je \(|\lambda| = 1\) in \(\lambda^n \neq\) za vsak \(n \in \NN\).
    \end{itemize}
\end{definicija}

\begin{definicija}
    Naj bo \(z^*\) fiksna točka preslikave \(f\). Množico
    \[\Omega (f, z^*) \coloneq \set{z_0 \in \CC : \lim_{n \to \infty} f^n (z_0) = z^*}\]
    imenujemo \emph{območje privlaka \(z^*\)}.
\end{definicija}

Množica \(A\) je za funkcijo \(f\) \emph{naprej invariantna}, če je \(f (A) = A\), in \emph{nazaj invariantna}, če je \(f^{-1} (A) = A\). Če je množica naprej in nazaj invariantna, pravimo, da je \emph{popolnoma invariantna}.

\begin{izrek}
    Naj bo \(f\) cela funkcija in \(z^*\) njena privlačna fiksna točka. Potem je \(\Omega (f, z^*)\) odprta in popolnoma invariantna.
\end{izrek}

\begin{dokaz}
    Naj bo \(a \in \Omega (f, z^*)\) in \(r > 0\), da je \(\overline{\Delta (z^*, 2r)} \subset \Omega (f, z^*)\). Ker \(f^n (a)\) konvergira k \(z^*\), obstaja \(k \in \NN\), da je \(|f^k (a) - z^*| < r\). Ker je \(f\) zvezna, obstaja \(\delta > 0\), da je \(|f^k (a) - f^k (z)| < r\) za vsak \(z \in \Delta (a, \delta)\). Potem je
    \[|f^k (z) - z^*| \leq |f^k (z) - f^k (a)| + |f^k (a) - z^*| \leq 2r\]
    za \(|z - a| < \delta\). Torej je \(f^k (z)\) vsebovana v \(\overline{\Delta (z^*, 2r)}\). Zaradi tega je \(\lim_{n \to \infty} f^{k + n} (z) = z^*\) za vsak \(z \in \Delta (a, \delta)\). Torej je \(\Delta (a, \delta) \subset \Omega (f, z^*)\).

    Naj bo \(z_0 \in \Omega (d, z^*)\). Potem
    \[\lim_{n \to \infty} f^n \prt{f (z_0)} = \lim_{n \to \infty} f^{n + 1} (z_0) = z^*\]
    in \(f (z_0) \in \Omega(f, z^*)\). Enako sklepamo, da \(f (z_0) \in \Omega (f, z^*)\) implicira \(z_0 \in \Omega (f, z^*)\).
\end{dokaz}

\noindent Komponenta za povezanost \(\Omega (f, z^*)\), ki vsebuje \(z^*\), se imenuje \emph{neposredno območje privlaka} \(z^*\). Označimo ga z \(\Omega_0 (f, z^*)\).

\begin{trditev} \label{prop:privlak}
    Naj bo \(f\) cela funkcija in \(\Omega \subsetneq \CC\) enostavno povezano območje, ki je popolnoma invariantno za \(f\). Če je \(z^{*} \in \Omega\) privlačna fiksna točka \(f\), potem je \(\Omega\) vsebovana v neposrednem območju privlaka \(z^*\).
\end{trditev}

Posvetimo se še nefiksnim periodičnim točkam. Če je \(z_0\) \(p\)-periodična, potem so vse točke \(z_0, f (z_0), f^2 (z_0), \dots, f^{p - 1} (z_0)\) fiksne točke preslikave \(f^p\). Pišemo \(z_j = f^j (z_0)\) in uporabimo verižno pravilo:
\[(f^p)' (z_j) = f' \prt{f^{p - 1} (z_j)} f' \prt{f^{p - 2} (z_j)} \dots f' \prt{f (z_j)} f' (z_j) = \prod_{k = 0}^{p - 1} f' (z_k).\]
To pomeni, da imajo vse točke \(z_j\) isti multiplikator \(K\). Če je \(K < 1\), potem rečemo, da je \(p\)-cikel privlačen. Če pa je \(K > 1\), rečemo, da je \(p\)-cikel odbojen.

\subsection{Normalne družine}

Naj bosta \(X\) in \(Y\) topološka prostora. Množico zveznih preslikav \(f \colon X \to Y\) opremimo s kompaktno-odprto topologijo. Tedaj relativno kompaktnim množicam pravimo \emph{normalne družine}. V tem poglavju jih bomo prek topologije enakomerne konvergence po kompaktih definirali za holomorfne funkcije na nekem območju \(D \subseteq \CC\). Vemo namreč, da sta zgornji topologiji na metričnih prostorih enaki.
\begin{definicija}
    Zaporedje holomorfnnih funkcij \((f_n)\) na območju \(D \subseteq \CC\) konvergira k \(f\) \emph{enakomerno po kompaktih}, če za vsako kompaktno podmnožico \(K \subset D\) in vsak \(\varepsilon > 0\) obstaja obstaja \(n_0 \in \NN\), da za \(n \geq n_0\) velja \(|f_n (z) - f(z)| < \varepsilon\) za vse \(z \in K\).
\end{definicija}

\begin{primer} \label{ex:primer}
    Zaporedje \(f_n \colon \DD \to \CC\), \(f_n (z) = z^n\) konvergira k \(f \equiv 0\) enakomerno po kompaktih. Res, naj bo \(K \subset \DD\) kompaktna. Tedaj obstaja \(r < 1\), da za vsak \(z \in K\) velja \(|z| < r\). Zato velja
    \[\abs{f_n (z) - f (z)} = \abs{z^n - 0} < r^n.\]
    Trditev sledi, saj je za vsak \(\varepsilon > 0\) tudi \(r^n < \varepsilon\) za dovolj velik \(n\).
\end{primer}

\noindent Naslednji izrek nam pove, da je konvergenca enakomerno po kompaktih v družini holomorfnih funkcij smiselna, saj ohranja holomorfnost.

\begin{izrek}
    Naj bo \((f_n)\) zaporedje holomorfnih funkcij na območju \(D \subseteq \CC\), ki konvergira k \(f\) enakomerno po kompaktih. Potem je tudi \(f\) holomorfna na \(D\).
\end{izrek}

\begin{dokaz}
    Naj bo \(\overline{\Delta (a, r)} \subseteq D\). Ker je \(\overline{\Delta (a, r)}\) kompaktna, zaporedje \(f_n\) tu enakomerno konvergira proti \(f\). Poleg tega vsaka funkcija \(f_n\) za vsak \(z \in \Delta (a, r)\) zadošča Cauchyjevi integralski formuli
    \[f_n (z) = \frac{1}{2 \pi i} \int_{\partial \Delta (a, r)} \frac{f_n (w)}{w - z} \, \dd w.\]
    Ker je konvergenca enakomerna, lahko zamenjamo limito in integral
    \[f (z) = \lim_{n \to \infty} f_n (z) = \lim_{n \to \infty} \frac{1}{2 \pi i} \int_{\partial \Delta (a, r)} \frac{f_n (w)}{w - z} \, \dd w = \frac{1}{2 \pi i} \int_{\partial \Delta (a, r)} \frac{f (w)}{w - z} \, \dd w.\]
    Funkcijo na desni lahko razvijemo v vrsto okoli \(z = a\) in zato je holomorfna na \(\Delta (a, r)\). Torej je taka tudi \(f\).
\end{dokaz}

\begin{definicija}
    Naj bo \(\FFF\) družina holomorfnih funkcij na območju \(D \subseteq \CC\). Pravimo, da je družina \emph{normalna}, če za vsako zaporedje funkcij \((f_n) \subset \FFF\) obstaja podzaporedje \((f_{n_k})\), ki konvergira enakomerno po kompaktih k neki \(f \in \FFF\) ali pa k \(f \equiv \infty\).
\end{definicija}

\subsection{Fatoujeva in Juliajeva množica}

Pri preučevanju dinamike funkcij, lahko domeno razdelimo na dva množici. Prvi, na kateri je obnašanje funkcije ``predvidljivo'', pravimo Fatoujeva množica, drugi, kjer je obnašanje kaotično, pa pravimo Juliajeva množica. Natančno ju definiramo sledeče.

Za celo funkcijo \(f\) je \emph{Fatoujeva množica} sestavljena iz vseh točk \(z_0\), za katere velja, da je \((f^n)\) normalna na neki okolici točke \(z_0\). Označimo jo z \(F (f)\). Njen komplement \(J (f) \coloneq \CC \setminus F (f)\) je \emph{Juliajeva množica}. Po definiciji vidimo, da je Fatoujeva množica vedno odprta, Juliajeva pa vedno zaprta.

\begin{primer}
    Po primeru \ref{ex:primer} zaporedje iteracij funkcije \(f (z) = z^2\) na \(\DD\) konvergira k \(f \equiv 0\) enakomerno po kompaktih. Podobno premislimo, da na \(\CC \setminus \overline{\DD}\) konvergira k \(f \equiv \infty\) enakomerno po kompaktih. Na \(\partial \DD\) zaporedje divergira. Zato je \(F (f) = \CC \setminus \partial \DD\) in \(J (f) = \partial \DD\).
\end{primer}

\begin{izrek}
    Naj bo \(B (f)\) množica vseh točk \(z \in \CC\) z omejeno orbito. Tedaj velja
    \[\partial (B (f)) \subset J (f).\]
\end{izrek}

\begin{dokaz}
    Naj bo \(z_0 \in \partial (B (f))\). Recimo, da \(z_0 \notin J (f)\). Potem je \((f^n)\) normalna na neki kompaktni okolici \(U\) točke \(z_0\). Ker je točka \(z_0\) robna, obstajata \(z_1, z_2 \in U\), tako da je orbita točke \(z_1\) omejena, orbita točke \(z_2\) pa ne. Torej ima \((f^n (z_n))\) divergentno podzaporedje \((f^{n_k} (z_2))_k\). Ker pa je \((f^n)\) normalna na \(U\), obstaja podzaporedje zaporedja \((f^{n_k} (z_2))_k\), ki enakomerno konvergira na \(U\). Ker je \((f^n (z_1))\) omejena in \((f^{n_k} (z_2))_k\) divergira k \(\infty\), smo prišli do protislovja.
\end{dokaz}

\noindent Če na tem mestu upoštevamo izrek \ref{thm:orbits}, hitro vidimo, da je za kompleksno eksponentno preslikavo \(J (f) = \CC\) in \(F (f) = \emptyset\). To pa ne drži tudi za poljubno celo funkcijo. Če za primer navedemo \(f (z) \coloneq (e^z - 1) / 4\), potem \((f^n)\) na \(\DD\) konvergira proti \(0\) enakomerno po kompaktih. Torej Fatoujeva množica vsebuje cel enotski disk.

Kot zanimivost omenimo, da za polinome velja enakost \(\partial (B (f)) = J (f)\). V tem primeru zato množici \(B (f)\) pravimo \emph{napolnjena Juliajeva množica}. Posledično je zgornji izrek uporaben za vizualizacijo Juliajeve množice. Na njem temelji na primer \emph{algoritem vzvratne iteracije}. Slike, ki nastanejo, pa so pogosto zanimive, kot vidimo na sliki \ref{fig:julia}. Tu je upodobljena Juliajeva množica funkcije \(z^2 + c\) za \(c \approx \complexnum{-0.5125+0.5213i}\). Z modro in zeleno barvo obarvana obnočja so Fatoujeva množica, medtem ko so rumeno obarvana območja Juliajeva množica.
\begin{figure}
    \centering
    \includegraphics[width=1\textwidth]{julia_set.pdf}
    \caption[Juliajeva množica]{Juliajeva množica funkcije \(z^2 + c\), koda prirejeno po \href{https://commons.wikimedia.org/w/index.php?curid=75993344}{Wikipedia Commons}.}
    \label{fig:julia}
\end{figure}
